\paragraph{Current formal application context (Section1.4)}
Based on palette-alternating tree codes, we devise a deterministic interactive coding scheme against adversarial errors that approaches capacity and thus matches the rate obtained by [13]. The advantage of our coding scheme is that given an explicit construction of palette-alternating tree codes, our scheme is efficient. We believe that the recent progress on tree code constructions [11,24] may eventually lead to constructions of palette-alternating tree codes. The coding scheme suggested in [13], on the other hand, relies on a certain counting argument, and it is not clear to us how to obtain an efficient scheme based on these ideas.
\begin{authorclaim}{Theorem6 (application context only)}
Let $\epsilon>0$. Assume there exists an explicit palette-alternating tree code with rate $1-\epsilon$ and distance $\delta=\Omega(\epsilon\cdot\log^{-1}(1/\epsilon))$ (which, computational aspects aside, we know exists by Theorem4). Then, there exists an efficient deterministic coding scheme against $\epsilon$-fraction of adversarial errors with rate $1-O(\sqrt{\epsilon\log(1/\epsilon)})$.
\end{authorclaim}
